\documentclass[class=article,doc={kv-usage}]{cusdoc}
\IfFontExistsTF{Consolas}
  {\newfontfamily\consolas{Consolas}[
    NFSSFamily     = consolas,
    UprightFont    = *,
    BoldFont       = *-Bold,
    ItalicFont     = *-Italic,
    BoldItalicFont = *-BoldItalic,
  ]}
  {\newfontfamily\consolas{cmun}[
    NFSSFamily     = consolas,
    Extension      = .otf,
    UprightFont    = *btl,
    BoldFont       = *tb,
    ItalicFont     = *bto,
    BoldItalicFont = *st,
    HyphenChar     = None]}
\DeclareTextFontCommand\textcon{\consolas}

\makeatletter
\pdfstringdefDisableCommands{\let\textcon\@firstofone}
\makeatletter

% \usepackage{minted}
% \newminted[typelua]{lua}{autogobble,fontfamily=consolas}

\usepackage[luarun]{cindexdoc}
\usepackage{cindexemoji}
\newcommand*\LUA{\texorpdfstring{\textcon{lua}}{lua}}

\enablecombinedlist

\begin{document}

\section*{目录}
\templatecombinedlist
  [section={space.left=2.5em, space.hang=2.5em, width.name=2.5em},
    subsection={space.left=4em, space.hang=2.5em, width.name=2.5em}]
  [columns=2,outer-sep=0pt,sep=20pt,ragged,
    adj,adj-outer={-\dimexpr\paperwidth-1.7cm*2-\textwidth}]
  {toc}{}

\section{\LUA 类型}\label{sec:lua-types}

\subsection{字符相关}

一般情况下，cindex 中保存数据的类型都能使用 \textcon{== nil} 来判断它是否为无效值。

\begin{luadoc}[global]{Char}
  \begin{luasyntax}
    Char(number) -> Char;
    Char(Char)   -> Char;
    Char(string) -> Char | nil;
  \end{luasyntax}
该类型的构造器。
返回一个 \luaref{Char}，这是 UTF-32 编码的 Unicode 字符。返回值可能并非有效的 Unicode 字符。
它与 \textcon{nil} 可进行相等性比较，如果为无效值，则它与 \textcon{nil} 判定为相等（但不可替代 \textcon{nil}）。

如果输入为字符串，则可能返回 \textcon{nil}。如果字符串的第一个字符为有效的 UTF-8 编码的字符，就返回这个字符。
\end{luadoc}

\begin{luadoc*}{Char}
# __eq, __lt, __le, __tostring
\luaref{Char} 实现的元方法。
\end{luadoc*}

\begin{luadoc*}{Char}
# is_valid
! Char.is_valid(self) -> boolean;
判断是否为有效的 Unicode 字符。
\end{luadoc*}

\begin{luarun}{}
local c = Char('汉');
print(c:is_valid());
local c = Char(0xDFFF);
print(c:is_valid(), c == nil);
\end{luarun}

\begin{luadoc*}{}
# from_number,to_number
! Char.from_number(number) -> Char;
! Char.to_number(self) -> number;
从数字转为 \luaref{Char} 或转为数字。
\end{luadoc*}

\begin{luadoc*}{}
# cmp
! Char.cmp(Char?, Char?) -> number | nil;
比较两个字符，如果前者小于后者，返回负数，如果大于后者，返回正数，如果相等，返回 0。如果二者之一为 \textcon{nil}，若是字符无效，返回 0，否则返回 \textcon{nil}。
\end{luadoc*}

\begin{luadoc*}{}
#<@@=Char> from_string
! Char.from_string(string) -> (Char | nil, boolean);
若字符串开头为无效的 UTF-8 序列，则返回 \textcon{nil}，此时第二个返回值总为 \textcon{false}。当输入为有效的 UTF-8 字符串时，第一个返回值总是一个有效的 \luaref{Char}，当字符串只有一个字符时，第二个返回值为 \textcon{true}。
\end{luadoc*}

\begin{luarun}{}
print(Char.from_string('汉'));
print(Char.from_string('汉字'));
print(Char.from_string('\xFF'));
\end{luarun}

\begin{luadoc*}{}
# chars
! Char.chars(string) -> function() -> Char | nil;
返回一个迭代器函数，这个函数依次返回 UTF-8 字符串的所有字符。
\end{luadoc*}

\begin{luarun}{}
for c in Char.chars('CJK汉字') do
    print(c);
end;
\end{luarun}

\begin{luadoc*}{}
# has_han_info, han_info
! Char.has_han_info(self) -> boolean;
! Char.han_info(self) -> CjkInfo | nil;
判断该字符是否有 CJK 相关的信息，或获取该信息。

Unicode 区块 CJK Unified Ideographs、CJK Extension A \~{} J、
CJK Compatibility Ideographs、CJK Compatibility Ideographs Supplement
中的所有字符必定有此信息，但“〇”（\texttt{U+3007}）不再此列。
\end{luadoc*}

\begin{luarun}{}
print(Char('L'):has_han_info());
print(Char('汉'):has_han_info());
print(Char('〇'):has_han_info()); -- U+3007
\end{luarun}

\begin{luadoc*}{}
# han_ordered_strokes
! Char.han_ordered_strokes(self) -> OrderedStrokes | nil;
返回汉字的笔顺。字形为中国大陆字形。
\end{luadoc*}

\begin{luarun}{}
print(Char('L'):han_ordered_strokes());
print(Char('汉'):han_ordered_strokes());
print(Char('〇'):han_ordered_strokes()); -- U + 3007
\end{luarun}

\begin{luadoc*}{}
# script
! Char.script(self) -> number;
返回字符所属的 Unicode Script 值，\textcon{number} 为内部句柄，可使用 \luaref[icu.table]{Script} 表查询该值代表的字符串。
\end{luadoc*}

\begin{luarun}{}
local icu_data = require("icu.table");
print(icu_data.Script[ Char('L'):script() ]);
print(icu_data.Script[ Char('汉'):script() ]);
print(icu_data.Script[ Char('〇'):script() ]);
\end{luarun}

\begin{luadoc*}{}
# general_category
! Char.general_category(self) -> number;
返回字符所属的 Unicode General Category 值，\textcon{number} 为内部句柄，可使用 \luaref[icu.table]{GeneralCategory} 表查询该值代表的字符串。
\end{luadoc*}

\begin{luadoc*}{}
# is_alphabetic, is_numeric, is_math
! Char.is_alphabetic(self) -> boolean;
判断字符是否为字母/数字/数学符号，范围涵盖所有 Unicode 字符。
\end{luadoc*}

\begin{luadoc*}{}
# is_ideographic
! Char.is_ideographic(self) -> boolean;
判断字符是否具有 Unicode Ideographic 属性。不单单只有 CJK 字符有此属性。
\end{luadoc*}

\begin{luadoc*}{}
# simple_fold
! Char.simple_fold(self) -> Char;
折叠字符的大小写。字符在折叠大小写之后可进行忽略大小写的比较。

若要对字符串进行忽略大小写的比较，可使用 \luaref[StrRef]{eq_ignore_case}。
\end{luadoc*}

\begin{luadoc}[global]{StrRef}
不可变的 UTF-8 编码的字符串。它并不持有这字符串。
\end{luadoc}

\begin{luadoc*}{}
# __eq, __lt, __le, __len, __tostring
\luaref{StrRef} 实现的元方法。它除了可以与另一个 \luaref{StrRef} 进行比较外，还能与 \textcon{string} 进行比较。当它与 \textcon{nil} 比较相等时，仅当它指向空指针时，被判定为与 \textcon{nil} 相等。
\end{luadoc*}

\begin{luadoc*}{}
# is_nil, is_empty
! StrRef.is_nil(self) -> boolean;
\luaref[StrRef]{is_nil} 判断它是否指向空指针，\luaref[StrRef]{is_empty} 当 \luaref[StrRef]{is_nil} 为真或其长度为 0 时返回 \textcon{true}。
\end{luadoc*}

\begin{luadoc*}{}
# from_ptr, unsafe_new
! StrRef.from_ptr(const StrRef*) -> StrRef | nil;
! StrRef.unsafe_new(StrRef | string) -> StrRef;
! StrRef.unsafe_new(const uint8_t*, size_t?) -> StrRef;
不安全的构建方法。当它指向的值被销毁后，它本身不一定被销毁，因此可能出现悬垂引用。
\end{luadoc*}

\begin{luadoc*}{}
# to_string_escaped
! StrRef.to_string_escaped(self, all: boolean) -> string;
对部分不可见字符转义后输出。如果第二个参数被视为真，则对所有字符都进行转义。每个字符的转义形式为 \verb|\u{HHHH..}|，使用大写的 16 进制。
\end{luadoc*}

\begin{luarun}{}
print(StrRef.unsafe_new('\x1F汉〇L'):to_string_escaped());
print(StrRef.unsafe_new('\x1F汉〇L'):to_string_escaped(true));
\end{luarun}

\begin{luadoc*}{}
# cmp
! StrRef.cmp(StrRef, StrRef) -> number;
! StrRef.cmp(StrRef, string) -> number;
! StrRef.cmp(string, StrRef) -> number;
比较两个字符串。若前者小于后者，返回负数，若等于，返回 0，若大于，返回正数。
\end{luadoc*}

\begin{luadoc*}{}
# eq_ignore_case, cmp_ignore_case
! StrRef.eq_ignore_case (StrRef | string, StrRef | string) -> boolean;
! StrRef.cmp_ignore_case(StrRef | string, StrRef | string) -> number;
忽略大小写进行比较。若前者小于后者，返回负数，若等于，返回 0，若大于，返回正数。

某些字符在折叠大小写之后可能是多个字符，因此依次对 \luaref[Char]{simple_fold} 的结果进行比较，其结果并不一定与 \luaref[StrRef]{eq_ignore_case} 返回的值相等。
\end{luadoc*}

\begin{luarun}{}
print(StrRef.eq_ignore_case('Ab汉', 'ab汉'));
print(StrRef.eq_ignore_case('aß', 'ASS'));
print(StrRef.eq_ignore_case('aß', 'AS'));
\end{luarun}

\begin{luadoc*}{}
# slice
! StrRef.slice(self, start: number, end_: number?) -> StrRef;
返回从 \textcon{start} 到 \textcon{end_}（不含）的部分，索引从 \textbf{0} 开始。
如果 \textcon{end_} 未给出或超过字符串的长度，则该值视为字符串的长度。索引以字节为准，而不是字符。

一般来讲，凡是 cindex 的类型，它们的方法使用到索引时，都是从 0 开始，若是需要数字的范围，则都是前闭后开的。
\end{luadoc*}

\begin{luarun}{}
print(StrRef.unsafe_new('CJK汉字'):slice(1, 6));
\end{luarun}

\begin{luadoc*}{}
# head
! StrRef.head(StrRef | string) -> number | nil;
返回字符串的第一个字符。
\end{luadoc*}

\begin{luarun}{}
print(StrRef.unsafe_new('汉字'):head());
\end{luarun}

\begin{luadoc*}{}
# chars
! StrRef.chars(self) -> function() -> number | nil;
返回一个迭代器函数，这个函数依次返回 UTF-8 字符串的所有字符。
\end{luadoc*}

\begin{luadoc*}{}
# lines
! StrRef.lines(self) -> function() -> StrRef | nil;
返回一个迭代器函数，这个函数依次返回 UTF-8 字符串的所有行。行结束符不会包含在结果字符串中。
\end{luadoc*}

\begin{luarun}{}
local s = "abc\n汉字\n什么";
for line in StrRef.unsafe_new(s):lines() do
    print(line);
end
\end{luarun}

\begin{luadoc*}{}
# split
! StrRef.split(self, pattern: nil | Char | StrRef | string)
!@            -> function() -> StrRef | nil;
返回一个迭代器函数，这个函数依次返回 UTF-8 字符串的子串。分割方式由 \textcon{pattern} 指定。

如果模式为 \textcon{nil}，则按连续的 ASCII 空白分割，此时，它\textbf{总不返回空串}；
否则按指定的字符或字符串分割。如果字符串以 \textcon{pattern} 结尾，那么最后一项就是这个 \textcon{pattern} 之前的部分子串。

如果输入为空串，则返回的第一个值就是 \textcon{nil}。
\end{luadoc*}

\begin{luarun}{}
local s = "::std::os::ffi::";
for part in StrRef.unsafe_new(s):split("::") do
    print(part);
end
\end{luarun}

\begin{luarun}{}
local s = "  std  os ffi \n ";
for part in StrRef.unsafe_new(s):split() do
    print(part);
end
\end{luarun}

\begin{luadoc*}{}
# graphemes
! StrRef.graphemes(self) -> function() -> StrRef | nil;
返回一个迭代器函数，这个函数依次返回每个 \href{https://www.unicode.org/reports/tr29/#Grapheme_Cluster_Boundaries}{Grapheme Cluster}。
\end{luadoc*}

\begin{luarun}{}
local s = "ABC汉字。";
for cluster in StrRef.unsafe_new(s):graphemes() do
    print(cluster);
end
\end{luarun}

\begin{luadoc*}{}
# unicode_words
! StrRef.unicode_words(self) -> function() -> StrRef | nil;
返回一个迭代器函数，这个函数依次返回每个单词。
\end{luadoc*}

\begin{luadoc*}{}
# split_word_bounds
! StrRef.split_word_bounds(self) -> function() -> StrRef | nil;
返回一个迭代器函数，这个函数返回一个根据 \href{https://www.unicode.org/reports/tr29/#Word_Boundaries}{Word Boundaries} 分隔开来的子串。
\end{luadoc*}

\begin{luarun}{}
local s = [[ABC: (“汉字”)。]];
for word in StrRef.unsafe_new(s):unicode_words() do
    print(word);
end
\end{luarun}

\begin{luarun}{}
local s = [[ABC: (“汉字”)。]];
for part in StrRef.unsafe_new(s):split_word_bounds() do
    print(part);
end
\end{luarun}

\begin{luadoc*}{}
# trim_ascii_spaces
! StrRef.trim_ascii_spaces() -> StrRef | nil;
! StrRef.trim_ascii_spaces(start: boolean) -> StrRef | nil;
! StrRef.trim_ascii_spaces(start: boolean, end: boolean) -> StrRef | nil;
返回删除字符串前后空格的结果。如果是有效的字符串，不会返回 \textcon{nil}。默认情况下，前后的空格都被删除。
\end{luadoc*}

\begin{luadoc*}{}
# is_emoji
! StrRef.is_emoji(self) -> boolean;
判断字符串是否为 Emoji 序列，有的 Emoji 是由多个 Unicode 字符组成。
\end{luadoc*}

\begin{luarun}{}
print(#"😂", StrRef.unsafe_new("😂"):is_emoji());
print(#"🇨🇳", StrRef.unsafe_new("🇨🇳"):is_emoji());
print(#"🤦🏻‍♀️", StrRef.unsafe_new("🤦🏻‍♀️"):is_emoji());
print(#"👩🏻‍❤️‍💋‍👨🏻", StrRef.unsafe_new("👩🏻‍❤️‍💋‍👨🏻"):is_emoji());
print(#"汉", StrRef.unsafe_new("汉"):is_emoji());
print(#"L", StrRef.unsafe_new("L"):is_emoji());
\end{luarun}

\begin{luadoc}[global]{CjkInfo}
保存 CJK 相关的信息。这个类型拥有其数据。一般通过 \luaref[Char]{han_info} 获取。
\end{luadoc}

\begin{luadoc*}{}
# __eq, __tostring
实现的元方法。
\end{luadoc*}

\begin{luadoc*}{}
# radical_from_kind
! CjkInfo.radical_from_kind(kind: StrRef | string) -> number | nil;
从康熙部首的种类返回其康熙部首的内部表示。种类是 \href{https://www.unicode.org/Public/UCD/latest/ucd/CJKRadicals.txt}{Unicode CJKRadicals.txt 数据}的第一列，一般是 1 \~{} 214 的数字，有的还可带 1 至 3 个西文单引号，如 \textcon{123}，\textcon{213''} 分别表示第 123 个康熙部首，以及第 213 个康熙部首的第 2 个变体。
\end{luadoc*}

\begin{luarun}{}
  print(true)
\end{luarun}

\begin{luadoc*}{}
# is_nil
! CjkInfo.is_nil(self) -> boolean;
判断是否为无效值。
\end{luadoc*}

\begin{luadoc*}{}
# contains_char
! CjkInfo.contains_char(Into<Char>) -> boolean;
判断字符是否有 CJK 信息。它与 \luaref[Char]{has_han_info} 完全一致。
\end{luadoc*}

\begin{luadoc*}{}
# kx_radical_char
! CjkInfo.kx_radical_char(self) -> Char;
返回康熙部首的具体字符，这字符位于 Unicode 的康熙部首相关区块。某些康熙部首的变体在康熙部首相关区块中没有对应值，因此，返回的字符可能无效。
\end{luadoc*}

\begin{luadoc*}{}
# kx_ideography_char
! CjkInfo.kx_ideography_char(self) -> Char;
返回康熙部首的具体字符，这字符位于 Unicode 的 CJK 表意字相关区块。所有康熙部首及其变体在表意字区块都有对应值，因此有效的 \luaref{CjkInfo} 总返回有效的字符。
\end{luadoc*}

\begin{luadoc*}{}
# kangxi_radical_number
! CjkInfo.kangxi_radical_number(self) -> number | nil;
返回康熙部首的号码，目前的返回值为 1 \~{} 214。
\end{luadoc*}

\begin{luadoc*}{}
# kangxi_radical_ideography_slot
! CjkInfo.kangxi_radical_ideography_slot(self)
!@                                      -> number | nil;
返回康熙部首在 CJK 表意字区块中的对应值的代码点。康熙部首的变体被映射为原始的 214 个部首。
\end{luadoc*}

\begin{luadoc*}{}
# kx_radical_kind
! CjkInfo.kx_radical_kind(self) -> StrRef | nil;
返回康熙部首的字符串表示。由数字和末尾的 0 \~ 3 个 \textcon{'} 组成。
\end{luadoc*}

\begin{luadoc*}{}
# ucd_additional_strokes
! CjkInfo.ucd_additional_strokes(self) -> number;
返回 Unicode 数据中该字符除部首外的笔画数，该值\textbf{可能为负数}。
\end{luadoc*}

\begin{luadoc*}{}
# ucd_total_strokes
! CjkInfo.ucd_total_strokes(self) -> number;
返回 Unicode 数据中该字符的总笔画数。这与所选字形无关，另见 \luaref[CjkInfo]{total_strokes}。
\end{luadoc*}

\begin{luadoc*}{}
# total_strokes
! CjkInfo.total_strokes(self) -> number;
返回字符的总笔画数。不同字形标准，总笔画数可能不同，该函数使用中国大陆字形标准。它的值不一定和 \luaref[CjkInfo]{ucd_total_strokes} 返回的一致。
\end{luadoc*}

\begin{luadoc*}{}
# original_radical
! CjkInfo.original_radical(self) -> number | nil;
如果是康熙部首的变体，则返回其原始的康熙部首的内部表示，若就是原始部首，则返回原始部首的内部表示。对于有效的内部值，总是不返回 \textcon{nil}。
\end{luadoc*}

\begin{luadoc*}{}
# mandarin
! CjkInfo.mandarin(self) -> nil | (StrRef, number);
如果该字符有拼音信息，第一个返回值是它最常见的拼音，第二个返回值是它的声调。
\end{luadoc*}

\begin{luadoc*}{}
# cantonese
! CjkInfo.cantonese(self) -> nil | (StrRef, number);
如果该字符有粤语拼音信息，第一个返回值是它最常见的拼音，第二个返回值是它的声调。
\end{luadoc*}

\begin{luadoc}[global]{Emoji}
  \begin{luasyntax}
    Emoji(string | StrRef) -> Emoji | nil;
  \end{luasyntax}
如果字符串是一个有效的 Emoji 序列，则返回这个 Emoji，Emoji 类型是 Owned。
\end{luadoc}

\begin{luadoc*}{}
# __tostring
把 emoji 转为字符串。
\end{luadoc*}

\begin{luadoc*}{}
# name
! Emoji.name(self) -> StrRef | nil;
返回 emoji 的名称。一般情况下该方法不会返回 \textcon{nil}。
\end{luadoc*}

\begin{luadoc*}{}
# group
! Emoji.group(self) -> number;
返回 emoji 组别的内部表示。可使用 \luaref[icu.table]{EmojiGroup} 获取组别的名字。
\end{luadoc*}

\begin{luadoc}[global]{OrderedStrokes}
汉字的笔画。
\end{luadoc}

\begin{luadoc*}{}
# __len, __eq, __lt, __gt, __ipairs
\luaref{OrderedStrokes} 实现的元方法。\luaref{__ipairs} 返回的索引号从 1 开始。
\end{luadoc*}

\begin{luadoc*}{}
# __tostring
转为字符串的形式。字符串的每一项都是 1 \~{} 5 的某个数字，按顺序分别代表：横、竖、撇、点（捺）、折（弯钩）。
\end{luadoc*}

\begin{luadoc*}{}
# __index, get
! OrderedStrokes[number] -> number;
! OrderedStrokes.get(n: number) -> number;
获取第 $n$ 画。$n$ 从 0 开始。数字为 1 \~{} 5 的某个数字，按顺序分别代表：横、竖、撇、点（捺）、折（弯钩）。
\end{luadoc*}


\subsection{索引条目}

\begin{luadoc}[global]{IndexEntryRef}
未合并的索引。
\end{luadoc}

\begin{luadoc*}{}
# __eq
\luaref{IndexEntryRef} 可以与 \textcon{nil} 进行比较，当其为无效值时，被视为与 \textcon{nil} 相等。
\end{luadoc*}

\begin{luadoc*}{}
# is_nil
! IndexEntryRef.is_nil(self) -> boolean;
判断是否为 \textcon{nil}。
\end{luadoc*}

\begin{luadoc*}{}
# to_handle
! IndexEntryRef.to_handle(self) -> number;
返回一个句柄。一般来讲，\luaref{IndexEntryRef} 是不能作为数字的索引的，可转为句柄。
\end{luadoc*}

\begin{luadoc*}{}
# levels_count
! IndexEntryRef.levels_count(self) -> number;
返回此条目的层级数。
\end{luadoc*}

\begin{luadoc*}{}
# key_n, keys
! IndexEntryRef.key_n(self, n: number) -> StrRef | nil;
! IndexEntryRef.keys(self) -> function() -> StrRef | nil;
返回该条目第 $n$ 个排序关键字，$n$ 从 0 开始。或返回一个迭代器函数。
\end{luadoc*}

\begin{luadoc*}{}
# as_table
! IndexEntryRef.as_table(self) -> table;
! table := {
!@  levels: { { key?: string, display: string }, .. },
!@  range?: "None" | "Open" | "Close",
!@  page_commands?: string,
!@  pages?: { { kind: string, value: number }, .. },
!@  pages_raw?: string,
! }
把条目表示为 \LUA 表。当前 \textcon{kind} 的值有：
\begin{itemize}
  \item\textcon{Arabic}（阿拉伯数字），
  \item\textcon{alpha}（小写拉丁字母组成的 26 进制数字），
  \item\textcon{Alpha}（大写拉丁字母组成的 26 进制数字），
  \item\textcon{roman}（小写罗马数字），
  \item\textcon{Roman}（大写罗马数字），
  \item\textcon{ZhDigits}（中文数串，由“零〇一二三四五六七八九”组成），
  \item\textcon{ZhNumber}（中文数字）。
\end{itemize}
\end{luadoc*}

\begin{luadoc*}{}
# is_same
! IndexEntryRef.is_same(lhs, rhs) -> boolean;
判断两个条目的排序关键字是否完全相等。
\end{luadoc*}

\begin{luadoc}[global]{IndexEntryArray}
存储索引条目的数组。
\end{luadoc}

\begin{luadoc*}{}
# __len, __ipairs
\luaref{IndexEntryArray} 实现的元方法。\textcon{__ipairs} 返回的索引数字从 1 开始。
\end{luadoc*}

\begin{luadoc*}{}
# is_nil, is_empty
! IndexEntryArray.is_nil(self) -> boolean;
! IndexEntryArray.is_empty(self) -> boolean;
判断数组是否为无效值，或其不含任何值。
\end{luadoc*}

\begin{luadoc*}{}
# __index, get
! IndexEntryArray[number] -> IndexEntryRef;
! IndexEntryArray.get(self, index: number) -> IndexEntryRef;
用数字索引数组，索引从 0 开始。若索引大于或等于数组的长度，\luaref{IndexEntryRef} 并非一个有效值，可使用 \textcon{== nil} 或 \luaref[IndexEntryRef]{is_nil} 来判断是否有效。
\end{luadoc*}

\begin{luadoc*}{}
# push
! IndexEntryArray.push(self, IndexEntryRefMut);
! IndexEntryArray.push(self, table);
在数组末尾附加一个条目。\textcon{table} 必须能转为合法的 \luaref{IndexEntryRefMut}，表的结构与 \luaref[IndexEntryRef]{as_table} 返回的一致。
\end{luadoc*}

\begin{luadoc*}{}
# extend_from_iterator
! IndexEntryArray.extend_from_iterator(self, init);
! init := function() ->
!@                function() -> IndexEntryRefMut | table | nil;
把迭代器返回的每一项附加到数组中。\textcon{init} 返回一个迭代器函数。
\end{luadoc*}

\begin{luadoc}[global]{IndexEntryRefMut}
  \begin{luasyntax}
    IndexEntryRefMut()                 -> IndexEntryRefMut;
    IndexEntryRefMut(table)            -> IndexEntryRefMut;
    IndexEntryRefMut(IndexEntryRef)    -> IndexEntryRefMut;
    IndexEntryRefMut(IndexEntryRefMut) -> IndexEntryRefMut;
  \end{luasyntax}
\luaref{IndexEntryRefMut} 的构造函数。它创建一个新的空索引条目，或从传入的参数中构造出一个新的索引条目。\textcon{table} 的结构和 \luaref[IndexEntryRef]{as_table} 返回的一致。
\end{luadoc}

\begin{luadoc*}{}
# push_level
! IndexEntryRefMut.push_level(self, key: nil, display: string | StrRef)
!@                           -> boolean;
! IndexEntryRefMut.push_level(self, level: string | StrRef)
!@                           -> boolean;
! IndexEntryRefMut.push_level(self, key: string, display: string)
!@                           -> boolean;
! IndexEntryRefMut.push_level(self, key: StrRef, display: StrRef)
!@                           -> boolean;
在条目中插入一个层级。
\end{luadoc*}

\begin{luadoc*}{}
# set_range
! IndexEntryRefMut.set_range(self, range?: StrRef | string) -> boolean;
设置本条目是否是页码范围的起始或结束位置。\textcon{range} 可选值为 \textcon{None}、\textcon{Open}、\textcon{Close}，如果不设置，则相当于 \textcon{None}。返回值表示是否成功设置。
\end{luadoc*}

\begin{luadoc*}{}
# set_page_commands
! IndexEntryRefMut.set_page_commands(
!@  self,
!@  page_commands?: StrRef | string
!@) -> boolean;
设置本条目页码要使用的命令。如果未给出或为 \textcon{nil}，相当于清除已经设置的命令。
\end{luadoc*}

\begin{luadoc*}{}
# push_page
! IndexEntryRefMut.push_page(
!@  self,
!@  kind: StrRef | string,
!@  value: Into<number>
!@) -> boolean;
向本条目中插入一个页码段。页码种类如 \luaref[IndexEntryRef]{as_table} 所列。

一个条目能同时使用多级页码，如 \texttt{MMMMI-2026-c-8}，就依次使用了 \textcon{Roman}，\textcon{Arabic}，\textcon{alpha}，\textcon{Arabic} 四级页码，但只算作一个页码。
\end{luadoc*}

\begin{luadoc*}{}
# set_pages_raw
! IndexEntryRefMut.set_pages_raw(self, pages_raw: StrRef | string)
!@                              -> boolean;
设置页码的显示效果。在写入到输出文件时，总是使用这个值，页码的具体数值只用于排序等情形，因此它可以不和具体的值挂钩。
\end{luadoc*}

\begin{luadoc*}{}
# as_table
! IndexEntryRefMut.as_table(self) -> table;
转为 \LUA 表，它本身的值保持不变。表的结构与 \luaref[IndexEntryRef]{as_table} 返回的一致。
\end{luadoc*}

\begin{luadoc}[global]{MergedEntryRef}
合并后的索引条目。
\end{luadoc}

\begin{luadoc*}{}
# __eq, is_nil, to_handle
! MergedEntryRef.is_nil(self) -> boolean;
! MergedEntryRef.to_handle(self) -> number;
\luaref[MergedEntryRef]{is_nil} 用于判断是否为无效值。\luaref[MergedEntryRef]{to_handle} 用于把条目转为内部句柄。
\end{luadoc*}

\begin{luadoc*}{}
# levels_count, level_n, levels
! MergedEntryRef.levels_count(self) -> number;
! MergedEntryRef.level_n(self, n: number) -> StrRef | nil;
! MergedEntryRef.levels(self)
!@                     -> function() -> StrRef | nil;
返回条目的层级数，或返回对应的层级。
\end{luadoc*}

\begin{luadoc*}{}
# key_n, keys
! MergedEntryRef.key_n(self, n: number) -> StrRef | nil;
! MergedEntryRef.keys(self) -> function() -> StrRef | nil;
返回条目的排序关键字。关键字的个数个层级数一致。
\end{luadoc*}

\begin{luadoc*}{}
# pages_count
! MergedEntryRef.pages_count(self) -> number;
返回合并后的索引条目包含的页码（或页码区间）数量。
\end{luadoc*}

\begin{luadoc*}{}
# page_command_n, page_commands
! MergedEntryRef.page_command_n(self, n: number) -> StrRef | false | nil;
! MergedEntryRef.page_commands(self) -> function() -> StrRef | false | nil;
返回合并后的索引条目的页码在输出时要使用的命令。如果页码没有设置命令，返回 \textcon{false}。
\end{luadoc*}

\begin{luadoc*}{}
# pages_raw_n, pages_raw
! MergedEntryRef.pages_raw_n(self, n: number)
!@                          -> nil | StrRef | (StrRef, StrRef);
! MergedEntryRef.pages_raws(self)
!@            -> function() -> nil | StrRef | (StrRef, StrRef);
返回页码要输出的原始字符串。如果是单个页码则返回一个 \luaref{StrRef}，如果是页码区间，则返回两个 \luaref{StrRef}。
\end{luadoc*}

\begin{luadoc*}{}
# page_span_n, page_spans
! MergedEntryRef.page_span_n(self, n: number) -> number;
! MergedEntryRef.page_spans(self) -> function() -> number | nil;
返回页码区间的长度。如果是无效页码，返回 0，如果是单个页码，返回 1，否则返回区间的长度。
\end{luadoc*}

\begin{luadoc*}{}
# max_same_key
! MergedEntryRef.max_same_key(MergedEntryRef, MergedEntryRef) -> number;
返回两个条目拥有的相同排序键的个数。
\end{luadoc*}

\begin{luadoc*}{}
# as_table
! MergedEntryRef.as_table(self) -> table;
! table := {
!@  levels: { { key?: string, display: string }, .. },
!@  pages?: { page_entry, .. },
! }
! page_entry := {
!@  kind: "single",
!@  page_command: nil | string,
!@  pages: { { kind: string, value: number }, .. },
!@  pages_raw: string,
! } | {
!@  kind: "range",
!@  page_command: nil | string,
!@  start: { { kind: string, value: number }, .. },
!@  start_raw: string,
!@  end: { { kind: string, value: number }, .. },
!@  end_raw: string,
! }
以 \LUA 表的形式表示该条目。页码的种类见 \luaref[IndexEntryRef]{as_table}。
\end{luadoc*}


\subsection{实用类型}

\begin{luadoc}[path=Logger]{error,warn,info,debug,trace}
  \begin{luasyntax}
    Logger.error(string | StrRef) -> boolean;
  \end{luasyntax}
在日志中写入对应级别的信息。
\end{luadoc}

\begin{luadoc}[global]{IstInputStyle}
  \begin{luasyntax}
    IstInputStyle: const table;
  \end{luasyntax}
索引的输入格式。
\end{luadoc}

\begin{luadoc*}{}
# is_nil
! IstInputStyle.is_nil(self) -> boolean;
判断是否不存在任何输入格式。一般情况下，输入格式都是存在的。
\end{luadoc*}

\begin{luadoc*}{}
# keyword, arg_open, arg_close, range_open, range_close, level, actual, encap, quote, escape, page_compositor, comment, separator
! IstInputStyle.keyword() -> StrRef | nil;
返回输入格式对应关键字的值。
\end{luadoc*}

\begin{luadoc}[global]{IstOutputStyle}
  \begin{luasyntax}
    IstOutputStyle: const table;
  \end{luasyntax}
索引的输出格式。
\end{luadoc}

\begin{luadoc*}{}
# is_nil
! IstOutputStyle.is_nil(self) -> boolean;
判断是否不存在任何输出格式。一般情况下，输出格式都是存在的。
\end{luadoc*}

\begin{luadoc*}{}
# preamble, postamble, group_skip, heading_prefix, heading_suffix, headings_flag, numhead_positive, numhead_negative, symhead_positive, symhead_negative, item_0, item_1, item_2, item_01, item_x1, item_12, item_x2, delim_0, delim_1, delim_2, delim_n, delim_r, delim_t, encap_prefix, encap_infix, encap_suffix, page_precedence, suffix_2p, suffix_3p, suffix_mp, stroke_prefix, stroke_suffix, radical_prefix, radical_suffix, radical_simplified_flag, radical_simplified_prefix, radical_simplified_delimiter, radical_simplified_suffix
! IstOutputStyle.preamble() -> StrRef | nil;
返回输出格式对应关键字的值。
\end{luadoc*}

\begin{luadoc}{CliOptions}
  \begin{luasyntax}
    CliOptions: const {
   @  c: boolean,
   @  q: boolean,
   @  r: boolean,
   @  strict: boolean,
   @  z?: string,
   @  extra_options?: { string, .. },
    };
  \end{luasyntax}
命令行选项。\textcon{extra_options} 是命令行参数 \textcon{--} 之后的内容。
\end{luadoc}

\begin{luadoc}[path=CIndex]{Util,Ideographic,Entry}
  \begin{luasyntax}
    CIndex.Util: const table;
    CIndex.Ideographic: const table;
    CIndex.Entry: const table;
  \end{luasyntax}
cindex 提供的实用模块。
\end{luadoc}

\begin{luadoc*}{CIndex.Util}
# is_ctype
! CIndex.Util.is_ctype(type, value) -> boolean;
判断某个值是否为某个 cindex 类型。无法判别所有 \LUA 原生类型。

例如，\textcon{CIndex.Ideographic.is_ctype(StrRef, s)} 即判断 \textcon{s} 是否为 \luaref{StrRef} 类型。
\end{luadoc*}

\begin{luadoc*}{CIndex.Util}
# table_to_json_string
! CIndex.Util.table_to_json_string(table) -> string;
把表转为 JSON 字符串。
\end{luadoc*}

\begin{luadoc*}{CIndex.Util}
# value_from_json_string
! CIndex.Util.value_from_json_string(string) -> ??;
把 JSON 字符串转为某个 \LUA 基本类型。
\end{luadoc*}

\begin{luadoc*}{CIndex.Ideographic}
# group_by_pinyin, group_by_bihua, group_by_bushou
! CIndex.Ideographic.group_by_pinyin(Char) -> string;
返回字符按拼音/笔画数/部首的组别，组别见 \luaref[cindex.table]{GroupOrder}。

如果字符不位于 CJK 表意字相关区块或没有对应信息，则返回 \textcon{Symbols}，即认为它是符号组。

除了 CJK 表意字相关区块的字符外，还识别\cref{tab:extra-chars} 补充的字符的信息。
\end{luadoc*}

\begin{table}[h]
  \centering
  \begin{tabular}{c >{\ttfamily}l c c l}
  \toprule
  字符 & {\normalfont 代码点} & 拼音 & 部首 & 笔顺 \\ \midrule
    〇 & U+3007 & \pinyin{ling2} & 乙（5） & 5 \\
  \bottomrule
  \end{tabular}
  \caption{补充的字符信息}\label{tab:extra-chars}
\end{table}

\begin{luadoc*}{CIndex.Entry}
# compare_by_default, compare_by_pinyin, compare_by_bihua, compare_by_bushou
! CIndex.Entry.compare_by_default(MergedEntryRef, MergedEntryRef)
!@                               -> boolean;
根据指定的比较算法比较两个条目，如果前者应当排在后者之前，则返回 \textcon{true}。

总体来说，cindex 首先进行忽略大小写的比较，若相等，再区分大小写进行比较，顺序按照 Unicode 的顺序。若汉字与非汉字进行比较，总是非汉字排在前面，而不考虑它们在 Unicode 中的顺序。

使用读音排序时，没有读音数据的字符（包括生僻字）一律排在有读音的字符之前，汉字按其
最常用读音比较，读音相同汉字的按 Unicode 编码排序。

使用笔画数和笔顺排序时，笔画数小的排在前面，笔画数相同的，按横、竖、撇、点（捺）、折的顺序
逐笔画比较，仍然相同的按 Unicode 编码排序；生僻汉字没有笔顺信息的，排在同笔画数有笔顺的字后面。

使用部首和除部首笔画数排序时，部首按康熙字典 214 部首顺序排列，部首和笔画数相同的按
Unicode 编码排序。
\end{luadoc*}


\subsection{PCRE2 正则表达式}

\href{https://pcre2project.github.io/pcre2/}{PCRE2} 是 C 语言实现正则表达式库。关于它的相关信息和语法，见 \url{https://pcre2project.github.io/pcre2/doc/}。

cindex 对其做了简单封装。

\begin{luadoc}[global]{Regex}
  \begin{luasyntax}
    Regex(string | StrRef | const uint8_t*, options: number) -> Regex;
  \end{luasyntax}
创建一个正则表达式。无效的模式直接报错。
\end{luadoc}

\begin{luadoc*}{}
# jit_compile
! Regex.jit_compile(self) -> true | (false, string);
尝试使用 JIT 编译。返回是否成功。如果为 \textcon{false}，还返回相关错误码。
\end{luadoc*}

\begin{luadoc*}{}
# is_match
! Regex.is_match(self, subject: StrRef | string) -> boolean;
检查 \textcon{subject} 是否匹配正则表达式。
\end{luadoc*}

\begin{luadoc*}{}
# find
! Regex.find(self, subject: StrRef, start_pos: number) -> StrRef | nil;
! Regex.find(self, subject: string, start_pos: number) -> string | nil;
从字符串的 \textcon{start_pos} 位置开始，查找匹配该的正则表达式的子串并返回。不论是 \luaref{StrRef} 还是 \textcon{string}，索引均从 0 开始。
\end{luadoc*}

\begin{luadoc*}{}
# captures
! Regex.captures(self, subject: StrRef, start_pos: number)
!@              -> nil | { [number] = StrRef, [string] = StrRef, .. };
! Regex.captures(self, subject: string, start_pos: number)
!@              -> nil | { [number] = string, [string] = string, .. };
查找并返回字符串的捕获组。第 0 个索引是匹配的子串。
\end{luadoc*}

\begin{luadoc*}{}
# gmatch
! Regex.gmatch(self, subject: StrRef, start_pos: number)
!@            -> nil | { [number] = StrRef, [string] = StrRef, .. };
! Regex.gmatch(self, subject: string, start_pos: number)
!@            -> nil | { [number] = string, [string] = string, .. };
返回一个迭代器函数，其返回值依次是匹配的捕获。
\end{luadoc*}

\begin{luadoc*}{}
# gsub
! Regex.gsub(
!@  self,
!@  subject: StrRef | string,
!@  replacement: StrRef | string,
!@  options?: number,
!@  once?: boolean;
! ) -> string;
以正则表达式 \textcon{replacement} 替换 \textcon{subject} 中匹配的内容。
\end{luadoc*}

\begin{luadoc*}{}
# Constant
! Regex.Constant: const { [string] = number, .. };
这个表保存一些 PCRE2 定义的常量。
\end{luadoc*}

\begin{luadoc*}{Regex.Constant}
# UNSET, ZERO_TERMINATED
! Regex.Constant.UNSET: number;
\end{luadoc*}

\begin{luadoc*}{Regex.Constant}
# ALLOW_EMPTY_CLASS, ALT_BSUX, AUTO_CALLOUT, CASELESS, DOLLAR_ENDONLY, DOTALL, DUPNAMES, EXTENDED, FIRSTLINE, MATCH_UNSET_BACKREF, MULTILINE, NEVER_UCP, NEVER_UTF, NO_AUTO_CAPTURE, NO_AUTO_POSSESS, NO_DOTSTAR_ANCHOR, NO_START_OPTIMIZE, UCP, UNGREEDY, UTF, NEVER_BACKSLASH_C, ALT_CIRCUMFLEX, ALT_VERBNAMES, USE_OFFSET_LIMIT, EXTENDED_MORE, LITERAL, MATCH_INVALID_UTF, ALT_EXTENDED_CLASS
! Regex.Constant.ALLOW_EMPTY_CLASS: number;
\end{luadoc*}

\begin{luadoc*}{Regex.Constant}
# NOTBOL, NOTEOL, NOTEMPTY, NOTEMPTY_ATSTART, PARTIAL_SOFT, PARTIAL_HARD, DFA_RESTART, DFA_SHORTEST, SUBSTITUTE_GLOBAL, SUBSTITUTE_EXTENDED, SUBSTITUTE_UNSET_EMPTY, SUBSTITUTE_UNKNOWN_UNSET, SUBSTITUTE_OVERFLOW_LENGTH, NO_JIT, COPY_MATCHED_SUBJECT, SUBSTITUTE_LITERAL, SUBSTITUTE_MATCHED, SUBSTITUTE_REPLACEMENT_ONLY
! Regex.Constant.NOTBOL: number;
\end{luadoc*}

\begin{luadoc*}{Regex.Constant}
# INFO_ALLOPTIONS, INFO_ARGOPTIONS, INFO_BACKREFMAX, INFO_BSR, INFO_CAPTURECOUNT, INFO_FIRSTCODEUNIT, INFO_FIRSTCODETYPE, INFO_FIRSTBITMAP, INFO_HASCRORLF, INFO_JCHANGED, INFO_JITSIZE, INFO_LASTCODEUNIT, INFO_LASTCODETYPE, INFO_MATCHEMPTY, INFO_MATCHLIMIT, INFO_MAXLOOKBEHIND, INFO_MINLENGTH, INFO_NAMECOUNT, INFO_NAMEENTRYSIZE, INFO_NAMETABLE, INFO_NEWLINE, INFO_DEPTHLIMIT, INFO_RECURSIONLIMIT, INFO_SIZE, INFO_HASBACKSLASHC, INFO_FRAMESIZE, INFO_HEAPLIMIT, INFO_EXTRAOPTIONS
! Regex.Constant.INFO_ALLOPTIONS: number;
\end{luadoc*}


\section{可加载的外部库}

cindex 对加载外部 \LUA 文件和库做了限制，目前只允许加载某些特定的库。

\subsection{\textcon{cindex.table}}

\begin{luadoc}[global]{cindex.table}
cindex 的相关表。可使用 \textcon{require("cindex.table")} 加载。
\end{luadoc}

\begin{luadoc*}{}
# GroupOrder
! cindex.table.GroupOrder: const {
!@  "nil" = 0,
!@  "Symbols" = ??,
!@  "Numbers" = ??,
!@  "A" = ??, .., "Z" = ??,
!@  "BiHua1" = ??, .., "BiHua64" = ??,
!@  "BuShou1" = ??, .., "BuShou214" = ??,
!@  [number] = string, ..
!};
cindex 默认的各个组别及其排序。\textcon{Symbols} 是符号分组，\textcon{Numbers} 是数字分组，\textcon{BiHua..} 是笔画数分组，\textcon{BuShou..} 是康熙部首分组。每个分组对应一个数字。也可通过数字反查对应的组别名。
\end{luadoc*}

\begin{luadoc*}{}
# BiHuaToGroup
! BiHuaToGroup: const { [1] = string, .., [64] = string };
根据笔画数返回组别名。
\end{luadoc*}

\begin{luadoc*}{}
# BuShouToGroup
! BuShouToGroup: const { [number] = string, .. }
根据康熙部首的代码点（包括康熙部首相关区块和表意字相关区块）值返回组别名。
\end{luadoc*}

\begin{luadoc*}{}
# BuShouToChars
! BuShouToChars: const {
!@  [1] = { number, .. }, .., [214] = { number, .. },
!@  ["BuShou1"] = { number, .. }, .., ["BuShou214"] = { number, .. },
! };
康熙部首对应的表意字字符，如果有变体，则依次存储在对应的表中。
\end{luadoc*}


\subsection{\textcon{icu.table}}

\begin{luadoc}[global]{icu.table}
\luaref{icu.table} 保存 Unicode 相关的信息。可使用 \textcon{require("icu.table")} 加载。
\end{luadoc}

\begin{luadoc*}{}
# Script
! Script = { [number] = string, [string] = number, .. }
\luaref[Char]{script} 返回的句柄可用于索引此表。也可通过字符串反查对应的句柄。
\end{luadoc*}

\begin{luadoc*}{}
# GeneralCategory
! GeneralCategory = { [number] = string, [string] = number, .. }
\luaref[Char]{general_category} 返回的句柄可用于索引此表。也可通过字符串反查对应的句柄。
\end{luadoc*}

%% @ 等于 \leavevmode
\begin{luadoc*}{}
# EmojiGroup
! EmojiGroup = { [number] = group, [string] = number, .. }
! group := | "Activities"
!@         | "Animals And Nature"
!@         | "Flags"
!@         | "Food And Drink"
!@         | "Objects"
!@         | "People And Body"
!@         | "Smileys And Emotion"
!@         | "Symbols"
!@         | "Travel And Places"
!@         | ??
Emoji 的分组。目前共有 9 个组别。
\end{luadoc*}

\begin{luadoc*}{}
# PropertyShortNames
! PropertyShortNames = { [string] = string, .. }
属性值对应的短名。
\end{luadoc*}


\subsection{\textcon{digest} 哈希算法}


\subsection{\textcon{LPeg}}

\href{https://www.inf.puc-rio.br/~roberto/lpeg/lpeg.html}{\textcon{LPeg}} 是基于 Parsing Expression Grammars 的模式匹配库。
cindex 内置了此库，可使用 \textcon{require("lpeg")} 加载。详细用法见 \url{https://www.inf.puc-rio.br/~roberto/lpeg/lpeg.html}。


\setuplayout{preset=balance}
\printindex[docusage,multi={ragged,sep=1em,outer-skip=0pt,cols=2}]

\end{document}